\chapter{Gaussian-filtered terms of an interaction}
\label{chap:gaussian_filtered_terms}

This chapter treats the Gaussian-filtered terms of a gapped Hamiltonian
$H=\sum_k h_k$ given by an interaction on a finite system of sites, after
M.~B. Hastings, \emph{Solving gapped Hamiltonians locally}. A term $W$ is
replaced by its Heisenberg evolution averaged against a centred Gaussian in
time; the average is nearly local, nearly preserves the ground space, and is
approximated by the same average under the terms near the support of $W$.
The counting results on graphs and lattices that these estimates use close the
chapter.

\section{Gaussian-filtered terms}
\label{sec:gaussian_filtered_lemma1}

\begin{definition}[Gaussian-filtered operator]
    \label{def:hastings_gaussian_filtered_term}
    \lean{ManyBody.gaussianFilteredTerm}
    \leanok
    For $q>0$ and a matrix $W$, the \emph{Gaussian-filtered operator} is
    \begin{align}
        \widetilde W & =\frac{\Delta E}{\sqrt{2\pi q}}\int_{-\infty}^{\infty}
        e^{iHt}We^{-iHt}\,e^{-(t\Delta E)^2/(2q)}\,dt,
        \label{eq:hastings_gaussian_filtered_term}
    \end{align}
    and $\widetilde W=W$ for $q=0$. Applied to the centered parts
    $W\in\{H_L,H_B,H_R\}$, (\ref{eq:hastings_gaussian_filtered_term}) gives
    $\tilde H_L^0,\tilde H_B^0,\tilde H_R^0$ of
    \cite[lines 783--793]{Hastings2007AreaLaw}.
\end{definition}

\begin{lemma}[Spectral properties of filtered operators]
    \label{thm:hastings_gaussian_filtered_spectral}
    \lean{ManyBody.inner_gaussianFilteredTerm, ManyBody.inner_gaussianFilteredTerm_ground, ManyBody.norm_gaussianFilteredTerm_le, ManyBody.gaussianFilteredTerm_sum, ManyBody.gaussianFilteredTerm_self, ManyBody.sum_gaussianFilteredTerm, ManyBody.inner_eq_zero_of_mem_groundSpace, ManyBody.norm_sub_starProjection_gaussianFilteredTerm_le}
    \leanok
    \uses{def:hastings_gaussian_filtered_term}
    If $Hv=Ev$ and $Hw=E'w$, then
    \begin{align}
        \langle v,\widetilde Ww\rangle
         & =e^{-q((E-E')/\Delta E)^2/2}\,\langle v,Ww\rangle,
        \label{eq:hastings_gaussian_filtered_matrix_element}
    \end{align}
    so $\langle v,\widetilde Ww\rangle=\langle v,Ww\rangle$ when $E=E'=E_0$.
    The map $W\mapsto\widetilde W$ is linear, satisfies
    $\|\widetilde W\|\le\|W\|$ and $\widetilde H=H$, and thus
    $\sum_i\widetilde H_i=H$ whenever $\sum_iH_i=H$
    \cite[line 793]{Hastings2007AreaLaw}. For a ground vector $\psi$ and the
    ground-space projection $P_{\mathrm{gs}}$,
    $\|(1-P_{\mathrm{gs}})\widetilde W\psi\|\le\|W\|e^{-q/2}\|\psi\|$.
\end{lemma}
\begin{proof}\leanok
    In the $H$-eigenbasis, $\langle v,e^{iHt}We^{-iHt}w\rangle
        =e^{i(E-E')t}\langle v,Ww\rangle$, whose Gaussian integral is
    (\ref{eq:hastings_gaussian_filtered_matrix_element}). Unitaries preserve
    $\|W\|$, and $e^{iHt}He^{-iHt}=H$. Since $|E-E_0|\ge\Delta E$ on the
    excited subspace, each excited coefficient of $\widetilde W\psi$ is scaled
    by at most $e^{-q/2}$.
\end{proof}

\begin{lemma}[Filtered operators without ground components]
    \label{thm:hastings_filtered_no_ground_component}
    \lean{ManyBody.norm_gaussianFilteredTerm_apply_le_commutator, ManyBody.norm_gaussianFilteredTerm_apply_le_commutator_of_isGapped}
    \leanok
    \uses{thm:hastings_gaussian_filtered_spectral}
    If $H\psi=E_0\psi$ and $P_{\mathrm{gs}}W\psi=0$, then
    \begin{align}
        \|\widetilde W\psi\| & \le\Delta E^{-1}e^{-q/2}\,\|[H,W]\psi\|.
        \label{eq:hastings_filtered_no_ground_component}
    \end{align}
    In particular, (\ref{eq:hastings_filtered_no_ground_component}) holds for a
    unique gapped ground state $\Psi_0$ when $\langle\Psi_0,W\Psi_0\rangle=0$,
    hence for the centered parts $H_L,H_B,H_R$
    \cite[lines 794--815]{Hastings2007AreaLaw}.
\end{lemma}
\begin{proof}\leanok
    For an $H$-eigenvector $v_E$ with eigenvalue $E$,
    $\langle v_E,[H,W]\psi\rangle=(E-E_0)\langle v_E,W\psi\rangle$. Since
    $\langle v_{E_0},W\psi\rangle=0$ and $E-E_0\ge\Delta E$ on the excited
    subspace, (\ref{eq:hastings_gaussian_filtered_matrix_element}) gives
    $|\langle v_E,\widetilde W\psi\rangle|
        \le\Delta E^{-1}e^{-q/2}|\langle v_E,[H,W]\psi\rangle|$.
\end{proof}

\begin{lemma}[Locality of filtered operators]
    \label{thm:hastings_filtered_locality}
    \lean{ManyBody.norm_map_gaussianFilteredTerm_le_mass, ManyBody.norm_commutator_gaussianFilteredTerm_le_mass, GaussianFilter.measure_compl_Icc_le_div_sq, ManyBody.norm_commutator_gaussianFilteredTerm_le}
    \leanok
    \uses{def:hastings_gaussian_filtered_term}
    If $q>0$, $T,\varepsilon\ge0$, and
    $\|[e^{iHt}We^{-iHt},O]\|\le\varepsilon$ for all $|t|\le T$, then
    \begin{align}
        \|[\widetilde W,O]\| & \le\varepsilon
        +2\|W\|\,\|O\|\,e^{-(T\Delta E)^2/(2q)}.
        \label{eq:hastings_filtered_locality}
    \end{align}
    For $q=0$, $\widetilde W=W$ and $\|[W,O]\|\le\varepsilon$ from $t=0$
    \cite[lines 816--828]{Hastings2007AreaLaw}.
\end{lemma}
\begin{proof}\leanok
    Split (\ref{eq:hastings_gaussian_filtered_term}) at $|t|=T$. On $[-T,T]$
    the integrand norm is at most $\varepsilon$; on $\mathbb R\setminus[-T,T]$
    it is at most $2\|W\|\,\|O\|$, and with $h=q/\Delta E^2$ the substitution
    $t=T+s$ and $(T+s)^2\ge T^2+s^2$ bound the tail mass by $e^{-T^2/(2h)}$.
\end{proof}

\begin{lemma}[Localization of filtered operators by partial trace]
    \label{thm:hastings_filtered_partial_trace}
    \lean{ManyBody.norm_sub_regionExpectation_gaussianFilteredTerm_le_mass, ManyBody.norm_sub_regionExpectation_gaussianFilteredTerm_le}
    \leanok
    \uses{thm:hastings_filtered_locality}
    Let $\mathbb E_K$ be the normalized partial trace over $K^c$ tensored with
    $1_{K^c}$. If $q>0$, $T,\varepsilon\ge0$, and
    $\|e^{iHt}We^{-iHt}-\mathbb E_K(e^{iHt}We^{-iHt})\|\le\varepsilon$ for all
    $|t|\le T$, then $\mathbb E_K(\widetilde W)$ is supported on $K$, satisfies
    $\|\mathbb E_K(\widetilde W)\|\le\|W\|$, and
    \begin{align}
        \|\widetilde W-\mathbb E_K(\widetilde W)\| & \le\varepsilon
        +2\|W\|\,e^{-(T\Delta E)^2/(2q)}.
        \label{eq:hastings_filtered_partial_trace}
    \end{align}
    For $q=0$, $\|\widetilde W-\mathbb E_K(\widetilde W)\|\le\varepsilon$. This
    replaces the truncated-evolution approximation of
    \cite[lines 816--828]{Hastings2007AreaLaw}.
\end{lemma}
\begin{proof}\leanok
    Since $\mathrm{id}-\mathbb E_K$ is continuous and linear with
    $\|(\mathrm{id}-\mathbb E_K)(M)\|\le2\|M\|$, splitting the time integral at
    $|t|=T$ as in Lemma~\ref{thm:hastings_filtered_locality} gives
    (\ref{eq:hastings_filtered_partial_trace}).
\end{proof}


\section{Commuting and zero-range terms}
\label{sec:gaussian_filtered_zero_range}

\begin{lemma}[Filtering an observable that commutes with the Hamiltonian]
    \label{thm:hastings_commuting_filter_exact}
    \lean{ManyBody.gaussianFilteredTerm_eq_of_commute, ManyBody.mem_groundSpace_apply_of_commute, ManyBody.starProjection_apply_eq_of_commute}
    \leanok
    \uses{def:hastings_gaussian_filtered_term}
    If $HW=WH$, then the Gaussian filter of $W$ under $H$ equals $W$ for
    every nonnegative variance. Moreover, if $H\psi=E_0\psi$, then
    $H(W\psi)=E_0W\psi$. Thus the orthogonal projection $P_0$ onto the
    $E_0$-eigenspace satisfies $P_0W\psi=W\psi$.
\end{lemma}
\begin{proof}\leanok
    Commutativity gives $e^{iHt}W=We^{iHt}$, so
    $e^{iHt}We^{-iHt}=W$. Its average against a Gaussian probability
    measure is $W$. Also $HW\psi=WH\psi=E_0W\psi$.
\end{proof}

\begin{lemma}[Commutativity of terms acting on individual sites]
    \label{thm:hastings_single_site_terms_commute}
    \lean{ManyBody.commute_of_singleton_support, ManyBody.commute_sum_of_singleton_support, ManyBody.gaussianFilteredTerm_sum_eq_of_singleton_support}
    \leanok
    Let $\Lambda$ be a finite set of sites, with a $d$-dimensional Hilbert
    space at each site. Suppose $H_i$ acts only at $i$ for every $i\in\Lambda$.
    Then $[H_i,H_j]=0$ for all $i,j$. Consequently, for every subset
    $s\subseteq\Lambda$, the sum $H_s=\sum_{j\in s}H_j$ commutes with every
    $H_i$, and the Gaussian filter of $H_i$ under $H_s$ equals $H_i$.
\end{lemma}
\begin{proof}\leanok
    \uses{thm:hastings_commuting_filter_exact}
    If $i\ne j$, the singleton supports are disjoint, so the terms commute.
    If $i=j$, the commutator also vanishes. Summing the commutation
    identities gives $[H_s,H_i]=0$, and
    Lemma~\ref{thm:hastings_commuting_filter_exact} evaluates the filter.
\end{proof}

\begin{lemma}[Exact truncation at zero interaction range]
    \label{thm:hastings_zero_range_exact}
    \lean{ManyBody.singleton_support_of_zero_range, ManyBody.gaussianFilteredTerm_sum_eq_of_zero_range, ManyBody.gaussianFilteredTerm_full_and_sum_eq_of_zero_range, ManyBody.starProjection_term_eq_of_zero_range, ManyBody.sub_starProjection_gaussianFilteredTerm_sum_eq_zero_of_zero_range}
    \leanok
    Let $\Lambda$ be a finite metric space, and suppose $H_i$ acts inside
    the closed ball of radius zero about $i$. Put $H=\sum_iH_i$ and let
    $s\subseteq\Lambda$ be arbitrary. For $q\ge0$ and $\Delta E\in\mathbb R$,
    both the full Gaussian filter $\widetilde H_i$ and the filter $M_i$
    under $H_s$ equal $H_i$. If $H\psi=E_0\psi$, then
    \begin{align*}
        \widetilde H_i & =M_i=H_i, \\
        (1-P_0)M_i\psi & =0.
    \end{align*}
    In particular, the approximation error in (apb) and the excited
    component in (facld) both vanish, without a spectral-gap assumption.
\end{lemma}
\begin{proof}\leanok
    \uses{thm:hastings_single_site_terms_commute, thm:hastings_commuting_filter_exact}
    A closed ball of radius zero in a metric space is the singleton of its
    centre. Apply Lemma~\ref{thm:hastings_single_site_terms_commute} to
    $s$ and $\Lambda$. The eigenspace statement in
    Lemma~\ref{thm:hastings_commuting_filter_exact} gives
    $P_0H_i\psi=H_i\psi$, hence the last identity.
\end{proof}


\section{Localization on a graph}
\label{sec:gaussian_filtered_local_projective_form}

\begin{lemma}[Locality of a filtered term on a graph]
    \label{thm:hastings05_filtered_locality_graph}
    \lean{ManyBody.norm_commutator_gaussianFilteredTerm_le_graphDistance}
    \leanok
    \uses{thm:hastings_filtered_locality, thm:dynamics_lieb_robinson}
    Let $W$ act on $S$, $B$ act on $\{y\}$ with $\min_{x\in S}d_G(x,y)\ge l$,
    and $q,\mu,T\ge0$, $\Delta E\in\mathbb R$. With $(C,v)$ from
    Theorem~\ref{thm:dynamics_lieb_robinson},
    \begin{align}
        \|[\widetilde W,B]\|\le C|S|\,\|W\|\,\|B\|\,e^{vT-\mu l}
        +2\|W\|\,\|B\|\,e^{-(T\Delta E)^2/(2q)}.
        \label{eq:hastings05_filtered_graph_commutator}
    \end{align}
    For $T=c_1l$ this is eq.~(lbnd) of~\cite[lines 266--283]{Hastings2006Solving}.
\end{lemma}
\begin{proof}\leanok
    Combine the $|t|\le T$ commutator bound of
    Theorem~\ref{thm:dynamics_lieb_robinson} with
    Lemma~\ref{thm:hastings_filtered_locality}.
\end{proof}

\begin{lemma}[The finite-range approximation of a filtered term]
    \label{thm:hastings05_filtered_approximation_graph}
    \lean{ManyBody.norm_gaussianFilteredTerm_sub_regionExpectation_le_graphNeighborhood, ManyBody.exists_local_approx_gaussianFilteredTerm, ManyBody.isHermitian_gaussianFilteredTerm}
    \leanok
    \uses{thm:hastings_filtered_partial_trace, thm:dynamics_localization, thm:sites_region_expectation, thm:hastings_gaussian_filtered_spectral}
    Assume $|\{y:d_G(x,y)=r\}|\le K_g(r+1)^k$ for all $x$ and $r\ge0$, and let
    $(C,v,T_k)$ be as in Lemma~\ref{thm:dynamics_localization}.
    For $W$ supported on $S$, $l\in\mathbb N$, $N_l(S)=\{y:d_G(y,S)\le l\}$,
    positive local dimension, $q,T\ge0$, and $\Delta E\in\mathbb R$, the
    operator $M=\mathbb E_{N_l(S)}(\widetilde W)$ acts on $N_l(S)$, satisfies
    $\|M\|\le\|W\|$ (and $M=M^\dagger$ if $W=W^\dagger$), and
    \begin{align}
        \|\widetilde W-M\|\le C|S|^2\,\|W\|\,K_gT_k\,e^{vT-l/2}
        +2\|W\|\,e^{-(T\Delta E)^2/(2q)}.
        \label{eq:hastings05_filtered_graph_approximation}
    \end{align}
    For $W=H_i$ supported in the $R$-ball around $i$ and
    $l=l_{\mathrm{proj}}-R\ge0$, $M_i=M$ acts within radius
    $l_{\mathrm{proj}}$ of $i$ with $\|M_i\|\le J$
    \cite[lines 285--320]{Hastings2006Solving},
    \cite{gap:hastings05_filtered_term_approximation}.
    Every finite graph satisfies this assumption with $K_g$ equal to its
    number of vertices and $k=0$, but that choice does not give a constant
    uniform in system size\cite{gap:hastings05_geometric_growth_scope}.
\end{lemma}
\begin{proof}\leanok
    Combine Lemma~\ref{thm:dynamics_localization} for $|t|\le T$
    with Lemma~\ref{thm:hastings_filtered_partial_trace}.
\end{proof}

\begin{lemma}[Ground vectors are nearly invariant under the approximation]
    \label{thm:hastings05_facld}
    \lean{ManyBody.norm_sub_starProjection_apply_le_of_gaussianFilteredTerm, ManyBody.norm_apply_sub_inner_smul_le_of_gaussianFilteredTerm}
    \leanok
    \uses{thm:hastings_gaussian_filtered_spectral, thm:hastings_filtered_no_ground_component, def:hamiltonians_gap}
    Let $H=H^\dagger$ have ground energy $E_0$, gap $\Delta E>0$, and ground
    projection $P$, and let $q\ge0$. For any $H\psi=E_0\psi$ and operators
    $W,M$, $\|(1-P)M\psi\|\le(\|W\|e^{-q/2}+\|\widetilde W-M\|)\|\psi\|$. If
    $\operatorname{ran}P=\mathbb C\Omega$ with $\|\Omega\|=1$, then
    \begin{align}
        \|M\Omega-\langle\Omega,M\Omega\rangle\Omega\|
        \le\|W\|e^{-q/2}+\|\widetilde W-M\|
        \label{eq:hastings05_filtered_ground_vector}
    \end{align}
    With $M=M_i$ as in
    Lemma~\ref{thm:hastings05_filtered_approximation_graph} and
    $q=c_1l_{\mathrm{proj}}\Delta E$, these give the form of eq.~(facld)
    and the unique-ground-state bound of
    \cite[lines 322--348 and 375--387]{Hastings2006Solving} under the
    additional growth assumption. They do not establish those statements
    on every source metric lattice\cite{gap:hastings05_geometric_growth_scope}.
\end{lemma}
\begin{proof}\leanok
    Since $(1-P)M\psi=(1-P)\widetilde W\psi+(1-P)(M-\widetilde W)\psi$,
    Lemma~\ref{thm:hastings_gaussian_filtered_spectral} bounds the first term
    by $\|W\|e^{-q/2}\|\psi\|$ and $\|(1-P)(M-\widetilde W)\psi\|\le\|\widetilde W-M\|\,\|\psi\|$.
    For $\operatorname{ran}P=\mathbb C\Omega$, $P\phi=\langle\Omega,\phi\rangle\Omega$.
\end{proof}

\begin{lemma}[The sum of the approximations]
    \label{thm:hastings05_sum_approximation}
    \lean{ManyBody.norm_sum_sub_le_of_sum_eq}
    \leanok
    \uses{thm:hastings_gaussian_filtered_spectral}
    If $H$ is Hermitian and $\sum_iW_i=H$, then any operators $M_i$ satisfy
    $\|\sum_iM_i-H\|\le\sum_i\|\widetilde W_i-M_i\|$. This is the comparison of
    $H$ with $M=\sum_iM_i$ in~\cite[lines 359--374]{Hastings2006Solving}. The
    consequence drawn there, that an actual ground state of $M$ has nonzero
    projection onto the ground eigenspace of $H$ at a radius of order
    $\log N$, is quantified in
    Lemma~\ref{thm:hastings05_metric_ground_overlap}.
\end{lemma}
\begin{proof}\leanok
    By Lemma~\ref{thm:hastings_gaussian_filtered_spectral},
    $\sum_i\widetilde W_i=H$, so $\sum_iM_i-H=\sum_i(M_i-\widetilde W_i)$.
\end{proof}


\section{Truncated dynamics}
\label{sec:gaussian_filtered_truncated_dynamics}

\begin{lemma}[Support of the truncated filter]
    \label{thm:hastings_truncated_support}
    \lean{ManyBody.hermitianUnitaryPath_conj_mem_supportedOperators, ManyBody.gaussianFilteredTerm_mem_supportedOperators}
    \leanok
    \uses{def:hastings_gaussian_filtered_term}
    If $H$ and $W$ act on $K$, then $e^{iHt}We^{-iHt}$ ($t\in\mathbb R$) and
    $\widetilde W$ act on $K$; in particular, the filter $M_i$ of $h_i$ under the terms within distance
    $l_{proj}-R$ of $i$ acts within distance $l_{proj}$ of $i$
    \cite[lines 302--304]{Hastings2006Solving}.
\end{lemma}
\begin{proof}
    \leanok
    Operators supported on $K$ form a closed subalgebra, hence contain
    $e^{iHt}We^{-iHt}$ and its Gaussian integral $\widetilde W$.
\end{proof}

\begin{lemma}[Filtering a comparison of evolutions]
    \label{thm:hastings_truncated_filter_comparison}
    \lean{ManyBody.norm_gaussianFilteredTerm_sub_le}
    \leanok
    \uses{def:hastings_gaussian_filtered_term}
    If $H,H'$ are Hermitian, $q,T,\varepsilon\ge0$, and
    $\|e^{iHt}We^{-iHt}-e^{iH't}We^{-iH't}\|\le\varepsilon$ for all $|t|\le T$,
    then the Gaussian filters $\widetilde W_H$ and $\widetilde W_{H'}$ satisfy
    $\|\widetilde W_H-\widetilde W_{H'}\|\le\varepsilon+2\|W\|e^{-(T\Delta E)^2/(2q)}$
    \cite[eq.~(dmi), lines 1524--1533]{Hastings2006Solving}.
\end{lemma}
\begin{proof}
    \leanok
    \uses{thm:hastings_filtered_locality}
    Split the Gaussian integral at $|t|=T$ as in
    Lemma~\ref{thm:hastings_filtered_locality}, using integrand bound
    $\varepsilon$ on $[-T,T]$ and $2\|W\|$ on $\mathbb R\setminus[-T,T]$.
\end{proof}

\begin{lemma}[Lieb--Robinson bound against an operator on several sites]
    \label{thm:hastings_truncated_multisite_lieb_robinson}
    \lean{ManyBody.norm_commutator_heisenbergEvolution_le_graphDistance_finset}
    \leanok
    \uses{thm:dynamics_lieb_robinson}
    Under the assumptions of Theorem~\ref{thm:dynamics_lieb_robinson}, if
    $A$ acts on $S$, $B$ acts on $Y$, and $d_G(x,y)\ge m$ for all $x\in S$,
    $y\in Y$, then for all $\mu\ge0$, with $C=2+4|b_0|e^{2\mu R}$ and
    $v=2v_R|b_0|e^{2\mu R}+1$,
    \begin{align}
        \|[B,\tau_t(A)]\|\le2|Y|\,C\,|S|\,\|A\|\,\|B\|\,e^{v|t|-\mu m}.
        \label{eq:hastings_truncated_multisite_lieb_robinson}
    \end{align}
\end{lemma}
\begin{proof}
    \leanok
    \uses{thm:dynamics_lieb_robinson, thm:sites_region_expectation}
    With $E$ the normalized partial trace over $Y$,
    $\|[B,\tau_t(A)]\|\le2\|B\|\,\|\tau_t(A)-E(\tau_t(A))\|$; bound the latter
    by summing the single-site commutators over $y\in Y$ via
    Lemma~\ref{thm:sites_region_expectation} and
    Theorem~\ref{thm:dynamics_lieb_robinson}.
\end{proof}

\begin{lemma}[The omitted terms against the truncated evolution]
    \label{thm:hastings_truncated_omitted_commutator}
    \lean{ManyBody.norm_commutator_sub_truncated_le}
    \leanok
    \uses{thm:hastings_truncated_multisite_lieb_robinson}
    Let $W$ act on $S\subseteq K$, $s$ be a family of terms supported in $K$,
    $H_s=\sum_{k\in s}h_k$, and $d_G(S,X_k)\ge m$ for every omitted term
    $k\notin s$ with $X_k\cap K\ne\emptyset$. Then
    \begin{align}
        \|[H-H_s,e^{iH_su}We^{-iH_su}]\|
        \le2v_RC\,|S|\,\|W\|\Bigl(\sum_{\substack{k\notin s \\ X_k\cap K\ne\emptyset}}\|h_k\|\Bigr)e^{v|u|-\mu m}
        \label{eq:hastings_truncated_omitted_commutator}
    \end{align}
    \cite[lines 1469--1505]{Hastings2006Solving}.
\end{lemma}
\begin{proof}
    \leanok
    \uses{thm:hastings_truncated_support, thm:hastings_truncated_multisite_lieb_robinson}
    By Lemma~\ref{thm:hastings_truncated_support}, $e^{iH_su}We^{-iH_su}$ acts
    on $K$, so omitted terms with $X_k\cap K=\emptyset$ commute with it; for
    $X_k\cap K\ne\emptyset$, apply
    (\ref{eq:hastings_truncated_multisite_lieb_robinson}) with $|X_k|\le v_R$.
\end{proof}

\begin{lemma}[Hastings' truncated approximation]
    \label{thm:hastings_truncated_approximation}
    \lean{ManyBody.norm_gaussianFilteredTerm_sub_truncated_le_sum, ManyBody.norm_gaussianFilteredTerm_sub_truncated_le, ManyBody.truncated_approx_gaussianFilteredTerm}
    \leanok
    \uses{def:hastings_gaussian_filtered_term}
    In the setting of Lemma~\ref{thm:hastings_truncated_omitted_commutator},
    if $\|h_k\|\le J$ and $N=|\{k\notin s:X_k\cap K\ne\emptyset\}|$, then the
    filter $M$ of $W$ under $H_s$ acts on $K$, satisfies $\|M\|\le\|W\|$ (and
    $M=M^\dagger$ if $W=W^\dagger$), and for all $q,T,\mu\ge0$,
    \begin{align}
        \|\widetilde W-M\|\le2v_RC\,|S|\,\|W\|\,J\,N\,T\,e^{vT-\mu m}
        +2\|W\|e^{-(T\Delta E)^2/(2q)}
        \label{eq:hastings_truncated_approximation}
    \end{align}
    \cite[eq.~(apb), lines 300--320]{Hastings2006Solving},
    \cite{gap:hastings05_truncated_dynamics}.
\end{lemma}
\begin{proof}
    \leanok
    \uses{thm:hastings_truncated_support, thm:dynamics_hamiltonian_comparison, thm:hastings_truncated_filter_comparison, thm:hastings_truncated_omitted_commutator, thm:hastings_gaussian_filtered_spectral}
    Combine Lemmas~\ref{thm:hastings_truncated_omitted_commutator},
    \ref{thm:dynamics_hamiltonian_comparison}, and
    \ref{thm:hastings_truncated_filter_comparison} with
    Lemmas~\ref{thm:hastings_truncated_support}
    and~\ref{thm:hastings_gaussian_filtered_spectral}.
\end{proof}

\begin{lemma}[The truncated approximation at the time window $T=q/\Delta E$]
    \label{thm:hastings_truncated_approximation_time}
    \lean{ManyBody.exp_neg_sq_div_eq_of_time_eq, ManyBody.norm_gaussianFilteredTerm_sub_truncated_le_of_time}
    \leanok
    \uses{thm:hastings_truncated_approximation}
    In the setting of Lemma~\ref{thm:hastings_truncated_approximation} with
    $q\ge0$ and $\Delta E>0$, setting $T=q/\Delta E$ gives
    \begin{align}
        \|\widetilde W-M\|\le2v_RC\,|S|\,\|W\|\,J\,N\,\frac{q}{\Delta E}\,
        e^{vq/\Delta E-\mu m}+2\|W\|e^{-q/2}
        \label{eq:hastings_truncated_approximation_time}
    \end{align}
    \cite[lines 337--349]{Hastings2006Solving}.
\end{lemma}
\begin{proof}
    \leanok
    \uses{thm:hastings_truncated_approximation}
    Substitute $T=q/\Delta E$ into (\ref{eq:hastings_truncated_approximation}).
\end{proof}

\begin{lemma}[Ground vectors under the truncated approximation]
    \label{thm:hastings_truncated_facld}
    \lean{ManyBody.norm_sub_starProjection_truncated_apply_le, ManyBody.norm_truncated_apply_sub_inner_smul_le}
    \leanok
    \uses{thm:hastings_truncated_approximation_time, thm:hastings05_facld}
    In the setting of Lemma~\ref{thm:hastings_truncated_approximation} with
    $q\ge0$, ground energy $E_0$, gap $\Delta E>0$, ground projection $P$, and
    \begin{align}
        \delta=3\|W\|e^{-q/2}+2v_RC\,|S|\,\|W\|\,J\,N\,\frac{q}{\Delta E}\,
        e^{vq/\Delta E-\mu m},
        \label{eq:hastings_truncated_facld_delta}
    \end{align}
    every ground vector $\psi$ satisfies $\|(1-P)M\psi\|\le\delta\|\psi\|$, and
    a unique unit ground state $\Omega$ satisfies
    $\|M\Omega-\langle\Omega,M\Omega\rangle\Omega\|\le\delta$
    \cite[eq.~(facld), lines 322--348 and 375--387]{Hastings2006Solving}.
\end{lemma}
\begin{proof}
    \leanok
    \uses{thm:hastings_truncated_approximation_time, thm:hastings05_facld}
    Combine Lemma~\ref{thm:hastings05_facld} with
    (\ref{eq:hastings_truncated_approximation_time}).
\end{proof}

\begin{lemma}[The sum of the truncated approximations]
    \label{thm:hastings_truncated_sum}
    \lean{ManyBody.norm_sum_truncated_sub_le, ManyBody.norm_sum_truncated_sub_le_card}
    \leanok
    \uses{thm:hastings_truncated_approximation_time, thm:hastings05_sum_approximation}
    In the setting of Lemma~\ref{thm:hastings_truncated_approximation} with
    $q,\mu\ge0$, $\Delta E>0$, and $H=\sum_ih_i=H^\dagger$, let each $h_i$
    (supported on $X_i\subseteq K_i$) be filtered under $\sum_{k\in s_i}h_k$ to
    $M_i$, where the $N_i$ omitted terms meeting $K_i$ satisfy
    $d_G(X_i,X_k)\ge m$. Then
    \begin{align}
        \Bigl\|\sum_iM_i-H\Bigr\|\le\sum_i\Bigl(2v_RC\,|X_i|\,\|h_i\|\,J\,N_i\,
        \frac{q}{\Delta E}\,e^{vq/\Delta E-\mu m}+2\|h_i\|e^{-q/2}\Bigr),
        \label{eq:hastings_truncated_sum}
    \end{align}
    which is at most $|\{i\}|\bigl(2v_R^2CJ^2N_0(q/\Delta E)e^{vq/\Delta E-\mu m}+2Je^{-q/2}\bigr)$
    when $N_i\le N_0$ \cite[lines 343--374]{Hastings2006Solving}.
\end{lemma}
\begin{proof}
    \leanok
    \uses{thm:hastings_truncated_approximation_time, thm:hastings05_sum_approximation}
    Combine Lemma~\ref{thm:hastings05_sum_approximation} with
    (\ref{eq:hastings_truncated_approximation_time}) for each $W=h_i$, and use
    $|X_i|\le v_R$, $\|h_i\|\le J$, $N_i\le N_0$.
\end{proof}

\begin{lemma}[Exponential decay under polynomial shell growth]
    \label{thm:hastings_truncated_facld_decay}
    \lean{ManyBody.exists_norm_sub_starProjection_truncated_apply_le_exp, ManyBody.exists_norm_truncated_apply_sub_inner_smul_le_exp, Real.add_one_pow_mul_le_mul_exp}
    \leanok
    \uses{thm:hastings_truncated_facld}
    In the setting of Lemma~\ref{thm:hastings_truncated_facld}, let $J\ge0$,
    fix $\mu>0$ and $c_1>0$ with $c_1v<\mu$, and set
    $l_{fac}=2/\min(c_1\Delta E,\ \mu-c_1v)$. For each projection radius
    $l\in\mathbb{N}$ let $M_l$ be the truncated filter of $W$ at
    $q=c_1l\Delta E$, built from the terms of a set $s_l$ acting inside a
    region $K_l\supseteq S$, such that every omitted term meeting $K_l$ lies at
    distance at least $m=l-3R$ from $S$. Suppose the number $N_l$ of omitted
    terms meeting $K_l$ satisfies $N_l\le K_g(l+1)^k$. Then there is a constant
    $C_1'\ge0$, independent of $l$, such that for every $l$ and every ground
    vector $\psi$ the component of $M_l\psi$ orthogonal to the ground space
    has norm at most $\|W\|(3+C_1'J)e^{-l/l_{fac}}\|\psi\|$.
    If $\Omega$ is a unique ground state with gap $\Delta E$, there is also a
    constant $C_2'\ge0$, independent of $l$, such that for every $l$,
    $\|M_l\Omega-\langle\Omega,M_l\Omega\rangle\Omega\|\le
        \|W\|(3+C_2'J)e^{-l/l_{fac}}$. This is the form $JO(e^{-l_{proj}/l_{fac}})$
    of eq.~(facld) of \cite[lines 322--349]{Hastings2006Solving}, with the
    polynomial growth of the shell as an additional hypothesis absent from
    source lines 169--190. Theorem~\ref{thm:hastings05_metric_facld} proves
    the actual physical metric cutoff without this extra growth condition
    \cite{gap:hastings05_truncated_dynamics,gap:hastings05_geometric_growth_scope}.
\end{lemma}
\begin{proof}
    \leanok
    \uses{thm:hastings_truncated_facld}
    With $\delta_0=\min(c_1\Delta E,\mu-c_1v)=2/l_{fac}>0$, $q=c_1l\Delta E$, and
    $\mu m\ge\mu l-3\mu R$, one has $e^{-q/2}\le e^{-\delta_0l/2}$,
    $e^{vq/\Delta E-\mu m}\le e^{3\mu R}e^{-\delta_0l}$, and
    $N_ll\le K_g(l+1)^kl\le K_g(k+1)!e^{\delta_0/2}(\delta_0/2)^{-(k+1)}e^{\delta_0l/2}$;
    substitute into (\ref{eq:hastings_truncated_facld_delta}).
\end{proof}

\begin{lemma}[Degree bounds and graph-ball separation]
    \label{thm:hastings05_degree_ball}
    \lean{ManyBody.card_graphBall_le, ManyBody.graphBall_subset_of_mem, ManyBody.graphBall_separation}
    \leanok
    \uses{def:dynamics_graph_ball}
    If every vertex has degree at most $D$, then
    $|B_G(x,r)|\le(D+1)^r$. Furthermore, $i\in B_G(x,l)$ implies
    $B_G(i,R)\subseteq B_G(x,l+R)$. If $k\notin B_G(x,l)$,
    $z\in B_G(x,R)$ and $y\in B_G(k,R)$, then
    $d_G(z,y)\ge\max(l-2R,0)$.
\end{lemma}
\begin{proof}\leanok
    \uses{def:dynamics_graph_ball}
    A shortest path shows that every vertex of $B_G(x,r+1)$ is either
    a vertex of $B_G(x,r)$ or a neighbour of one. Hence its cardinality
    is at most $(D+1)|B_G(x,r)|$, and induction proves the bound.
    The remaining assertions follow from the graph triangle inequality;
    for separation, the reverse inequality would place $k$ in $B_G(x,l)$.
\end{proof}

\begin{lemma}[Weighted local budgets]
    \label{thm:hastings05_weighted_region_budget}
    \lean{ManyBody.sum_weight_mul_comp_le_budget_mul_sum, ManyBody.sum_weight_meeting_le_card_mul_budget, ManyBody.sum_weight_meeting_graphBall_le}
    \leanok
    \uses{thm:hastings05_degree_ball}
    Let $w_k\ge0$ be weights on finitely many supports $X_k$ and assume
    $\sum_{k:y\in X_k}w_k\le B$ at every site $y$. For any family $T$
    of terms whose supports meet a region $K$,
    \begin{align}
        \sum_{k\in T}w_k & \le B|K|.
        \label{eq:hastings05_weighted_region_budget}
    \end{align}
    In particular, if the graph degree is at most $D$ and
    $K\subseteq B_G(x,l+C)$, then this sum is at most
    $B(D+1)^C e^{\nu l}$, where $\nu=\ln(D+1)$.
    More generally, for any anchors $a_k$, any real weights $w_k$,
    and $f\ge0$, the bound $\sum_{k:a_k=y}w_k\le B$ implies
    $\sum_kw_kf(a_k)\le B\sum_yf(y)$.
\end{lemma}
\begin{proof}\leanok
    \uses{thm:hastings05_degree_ball}
    Each term meeting $K$ contributes its weight to at least one site
    budget in $K$. Summing those budgets proves
    (\ref{eq:hastings05_weighted_region_budget}); apply the ball bound
    for its exponential form. The last assertion follows by grouping
    the sum according to the anchor and applying the fibre bound.
\end{proof}

\begin{lemma}[Exponentially weighted graph tails]
    \label{thm:hastings05_degree_weighted_tail}
    \lean{ManyBody.sum_exp_neg_edist_le_degree, ManyBody.sum_weight_exp_neg_edist_tail_le_degree}
    \leanok
    \uses{thm:hastings05_degree_ball, thm:hastings05_weighted_region_budget}
    Suppose the degree is at most $D$. Put $\nu=\ln(D+1)$, let
    $\mu>\nu$, and define $k_\mu(x,y)=e^{-\mu d_G(x,y)}$ at finite
    distance and $k_\mu(x,y)=0$ at infinite distance. Then
    \begin{align}
        \sum_y k_\mu(x,y) & \le\frac1{1-e^{-(\mu-\nu)}}.
        \label{eq:hastings05_degree_exponential_row}
    \end{align}
    If $w_k\ge0$, $B\ge0$ and $\sum_{k:a_k=y}w_k\le B$,
    then every family $T$ with $d_G(x,a_k)\ge m$ for $k\in T$ satisfies
    \begin{align}
        \sum_{k\in T}w_k k_\mu(x,a_k)
         & \le\frac{B e^{-\beta m}}{1-e^{-\beta}},
        \qquad \beta=\frac{\mu-\nu}{2}>0.
        \label{eq:hastings05_degree_exponential_tail}
    \end{align}
    Thus summation over omitted terms requires only bounded local geometry
    when the spatial decay rate is chosen above the graph growth rate.
\end{lemma}
\begin{proof}\leanok
    \uses{thm:hastings05_degree_ball, thm:hastings05_weighted_region_budget}
    Group vertices by their finite distance $r$ from $x$. Each sphere has
    at most $(D+1)^r$ vertices, so the sum is bounded by the geometric
    series $\sum_{r\ge0}e^{-(\mu-\nu)r}$, proving
    (\ref{eq:hastings05_degree_exponential_row}). On the omitted family,
    $k_\mu(x,a_k)\le e^{-\beta m}k_{\mu-\beta}(x,a_k)$.
    Grouping by anchors and applying
    (\ref{eq:hastings05_degree_exponential_row}) at rate $\mu-\beta$
    proves (\ref{eq:hastings05_degree_exponential_tail}).
\end{proof}

\begin{lemma}[Absorbing the time prefactor]
    \label{lem:hastings05_gaussian_decay_absorption}
    \leanok
    Let $a,\gamma>0$ and $A,w,l\ge0$, and put $\delta=\min(a,\gamma)$.
    Then
    \begin{align}
        3we^{-al/2}+Awl e^{-\gamma l}
        \le w\left(3+\frac{A}{e\delta/2}\right)e^{-\delta l/2}.
        \label{eq:hastings05_gaussian_decay_absorption}
    \end{align}
    In the truncated Gaussian approximation, $a=c_1\Delta E$ and
    $\gamma$ is the spatial decay rate after subtracting $c_1v$.
\end{lemma}
\begin{proof}\leanok
    The Gaussian term is at most $3we^{-\delta l/2}$. For the second term,
    $le^{-\gamma l}\le le^{-\delta l}$ and
    $le^{-\delta l}\le(e\delta/2)^{-1}e^{-\delta l/2}$, by maximizing
    $u e^{-u}$ at $u=1$ with $u=\delta l/2$.
\end{proof}

\begin{lemma}[The spectral gap above the actual ground eigenspace]
    \label{thm:hastings05_ground_space_quadratic_gap}
    \lean{ManyBody.groundSpace_gap_le_of_eigenvalues}
    \leanok
    Let $H$ be Hermitian and suppose every eigenvalue equals $E_0$ or is
    at least $E_0+\Delta$. Let $P_0$ be the orthogonal projection onto
    $\ker(H-E_0)$. Then every vector $\psi$ satisfies
    \begin{align}
        \Delta\bigl(\|\psi\|^2-\|P_0\psi\|^2\bigr)
        \le \operatorname{Re}\langle\psi,H\psi\rangle-E_0\|\psi\|^2.
        \label{eq:hastings05_ground_space_quadratic_gap}
    \end{align}
\end{lemma}
\begin{proof}\leanok
    In an orthonormal eigenbasis $b_k$, the excess energy is
    $\sum_k(\lambda_k-E_0)|\langle b_k,\psi\rangle|^2$.
    The coefficients of $(1-P_0)\psi$ vanish when $\lambda_k=E_0$ and
    otherwise equal those of $\psi$. Each remaining coefficient has
    weight at least $\Delta$. Parseval and the orthogonal decomposition
    $\psi=P_0\psi+(1-P_0)\psi$ give the inequality.
\end{proof}

\begin{lemma}[An actual ground vector of the approximate Hamiltonian]
    \label{thm:hastings05_approximate_ground_vector_energy}
    \lean{ManyBody.exists_groundVector_energy_le_of_norm_sub_le}
    \leanok
    Let $M$ be Hermitian, $\|H-M\|\le\epsilon$, and let $\Omega$ be a unit
    eigenvector of $H$ of real energy $E_0$. There exist a real number $e$
    and a unit vector $\phi$ such that $M\phi=e\phi$,
    $e\|\psi\|^2\le\operatorname{Re}\langle\psi,M\psi\rangle$ for every
    $\psi$, and
    \begin{align}
        \operatorname{Re}\langle\phi,H\phi\rangle-E_0\le2\epsilon.
        \label{eq:hastings05_approximate_ground_vector_energy}
    \end{align}
    No spectral-gap hypothesis on $M$ is needed.
\end{lemma}
\begin{proof}\leanok
    Minimize the Rayleigh quotient of $M$ on the unit sphere. A minimizer
    is an eigenvector. The variational principle and the two norm-error
    estimates give
    $\langle\phi,H\phi\rangle
        \le\langle\phi,M\phi\rangle+\epsilon
        \le\langle\Omega,M\Omega\rangle+\epsilon
        \le E_0+2\epsilon$, after taking real parts.
\end{proof}

\begin{lemma}[Ground-space overlap and phase accuracy]
    \label{thm:hastings05_ground_overlap_norm_error}
    \lean{ManyBody.exists_groundVector_projected_mass_of_norm_sub_le, ManyBody.exists_groundVector_overlap_of_norm_sub_le}
    \leanok
    \uses{thm:hastings05_approximate_ground_vector_energy, thm:hastings05_ground_space_quadratic_gap}
    Under the quadratic gap (\ref{eq:hastings05_ground_space_quadratic_gap})
    and the hypotheses of
    Lemma~\ref{thm:hastings05_approximate_ground_vector_energy}, an actual
    unit ground vector of $M$ can be chosen with
    \begin{align}
        \Delta(1-\|P_0\phi\|^2)\le2\epsilon.
        \label{eq:hastings05_ground_overlap_norm_error}
    \end{align}
    Consequently $P_0\phi\ne0$ whenever $2\epsilon<\Delta$.
    If the original ground state is unique, with unit vector $\Omega$
    and gap $\Delta>0$, a unit ground vector of $M$ also satisfies
    \begin{align}
        \Delta(1-|\langle\Omega,\phi\rangle|^2) & \le2\epsilon,        \\
        \min_{\theta\in\mathbb R}
        \|\phi-e^{i\theta}\Omega\|^2            & \le4\epsilon/\Delta.
        \label{eq:hastings05_ground_phase_norm_error}
    \end{align}
    The minimum is attained. These are quantitative forms of the overlap
    assertion of~\cite[lines 362--367]{Hastings2006Solving}.
\end{lemma}
\begin{proof}\leanok
    \uses{thm:hastings05_approximate_ground_vector_energy, thm:hastings05_ground_space_quadratic_gap}
    Apply the gap inequality to the minimizing vector. In the unique case,
    $P_0\phi=\langle\Omega,\phi\rangle\Omega$ and the minimum squared
    phase distance is $2(1-|\langle\Omega,\phi\rangle|)$, which is at most
    $2(1-|\langle\Omega,\phi\rangle|^2)$.
\end{proof}

\begin{lemma}[The logarithmic overlap radius]
    \label{thm:hastings05_logarithmic_overlap_radius}
    \lean{ManyBody.exists_groundVector_projection_of_exp_norm_bound, ManyBody.exists_groundVector_overlap_of_exp_norm_bound}
    \leanok
    \uses{thm:hastings05_ground_overlap_norm_error}
    Let $\beta,\Delta,\eta>0$, and suppose the Hermitian approximation
    satisfies $\|H-M_\ell\|\le B e^{-\beta\ell}$. Put
    $\tau=\min(\eta,1)$ and $Q=\max(B,1)$. If
    \begin{align}
        \ell\ge\frac1\beta\ln\frac{4Q}{\Delta\tau},
        \label{eq:hastings05_logarithmic_overlap_radius}
    \end{align}
    then $B e^{-\beta\ell}\le\Delta\tau/4$. Under the original ground-space
    gap, there is an actual unit ground vector $\phi$ of $M_\ell$ with
    $P_0\phi\ne0$ and $1-\|P_0\phi\|^2\le\eta$. In the unique-ground-state
    case there is an actual unit ground vector of nonzero overlap with
    $\Omega$ whose attained minimum squared phase error is at most $\eta$.
\end{lemma}
\begin{proof}\leanok
    \uses{thm:hastings05_ground_overlap_norm_error}
    Exponentiating (\ref{eq:hastings05_logarithmic_overlap_radius}) gives
    $e^{-\beta\ell}\le\Delta\tau/(4Q)$. Since $B\le Q$, the norm error is
    at most $\Delta\tau/4$. Apply
    Lemma~\ref{thm:hastings05_ground_overlap_norm_error}; its projection
    cannot vanish because $2\epsilon\le\Delta/2<\Delta$.
\end{proof}

\begin{lemma}[Ground-state overlap for the actual metric truncation]
    \label{thm:hastings05_metric_ground_overlap}
    \lean{ManyBody.exists_metric_truncated_groundVector_projection, ManyBody.exists_metric_truncated_groundVector_overlap}
    \leanok
    \uses{thm:hastings05_metric_filter_sum, thm:hastings05_ground_space_quadratic_gap, thm:hastings05_logarithmic_overlap_radius}
    Fix the local parameters $R>0,S,J\ge0,\Delta>0$. There are
    $c_1,\beta>0$ and $C\ge0$, independent of the finite metric space and
    Hamiltonian, with the following property. Let $H=\sum_i h_i$ satisfy
    the support, norm and local-cardinality assumptions of
    Theorem~\ref{thm:hastings05_metric_filter_sum}, and suppose its spectrum
    consists of $E_0$ and eigenvalues at least $E_0+\Delta$. Assume a unit
    ground vector exists, and let $N$ be the number of sites. For every
    $\eta>0$ and every
    \begin{align}
        \ell\ge\max\left\{R,\frac1\beta
        \ln\frac{4\max(NJC,1)}{\Delta\min(\eta,1)}\right\},
        \label{eq:hastings05_metric_ground_overlap_radius}
    \end{align}
    the actual truncated sum
    $M_\ell=\sum_i M_i(c_1\ell\Delta,\Delta,\ell)$ has a unit ground vector
    $\phi$ with nonzero projection onto the actual ground eigenspace of
    $H$, and $1-\|P_0\phi\|^2\le\eta$. If the original ground state is
    unique, an actual unit ground vector of $M_\ell$ has nonzero overlap
    with it and attained minimum squared phase error at most $\eta$.
    For fixed local parameters and accuracy, the required radius grows
    as $O(\log N)$, as asserted in
    \cite[lines 362--367]{Hastings2006Solving}.
\end{lemma}
\begin{proof}\leanok
    \uses{thm:hastings05_metric_filter_sum, thm:hastings05_ground_space_quadratic_gap, thm:hastings05_logarithmic_overlap_radius}
    Each truncated term, hence their sum, is Hermitian.
    Theorem~\ref{thm:hastings05_metric_filter_sum} bounds the norm error by
    $NJC e^{-\beta\ell}$. The source's eigenvalue gap gives the quadratic
    gap above the canonical ground eigenspace by
    Lemma~\ref{thm:hastings05_ground_space_quadratic_gap}. Apply
    Lemma~\ref{thm:hastings05_logarithmic_overlap_radius} with $B=NJC$.
\end{proof}


\section{Truncation in a finite metric space}
\label{sec:gaussian_filtered_metric_localization}

\begin{lemma}[Termwise commutator control of Gaussian truncation]
    \label{thm:hastings05_termwise_gaussian_truncation}
    \lean{ManyBody.norm_gaussianFilteredTerm_sub_truncated_le_of_commutator_bounds}
    \leanok
    Let $H=\sum_jh_j$, with every $h_j$ Hermitian. Suppose $C_j\ge0$ and
    $\|[h_j,\tau_t^s(W)]\|\le C_j$ whenever $j\notin s$ and $|t|\le T$.
    Then
    \begin{align}
        \|\mathcal F_{q,\Delta}^{H}(W)-\mathcal F_{q,\Delta}^{H_s}(W)\|
         & \le T\sum_{j\notin s}C_j
        +2\|W\|e^{-(T\Delta)^2/(2q)}.
        \label{eq:hastings05_termwise_gaussian_truncation}
    \end{align}
\end{lemma}
\begin{proof}\leanok
    \uses{thm:dynamics_hamiltonian_comparison, thm:hastings_truncated_filter_comparison}
    The finite-sum triangle inequality gives
    $\|[H-H_s,\tau_t^s(W)]\|\le\sum_{j\notin s}C_j$.
    The Duhamel formula therefore bounds the difference of the two
    evolutions by $T\sum_{j\notin s}C_j$ on $[-T,T]$.
    Gaussian integration adds the tail in
    (\ref{eq:hastings05_termwise_gaussian_truncation}).
\end{proof}

\begin{lemma}[Propagation between anchored supports]
    \label{thm:hastings05_anchor_lr}
    \lean{ManyBody.norm_commutator_heisenbergEvolution_le_anchor_kernel}
    \leanok
    Suppose $B$ acts on $Y$, with $d_G(y,b)\le r$ for $y\in Y$.
    For every $\mu>0$ and $t\in\mathbb R$,
    \begin{align}
        \|[B,e^{iHt}We^{-iHt}]\|
         & \le A_\mu(S)\|W\|\|B\|e^{v_\mu|t|}k_\mu(z,b).
        \label{eq:hastings05_anchor_lr}
    \end{align}
\end{lemma}
\begin{proof}\leanok
    \uses{thm:dynamics_lieb_robinson}
    If $d_G(z,b)>2r+l$, the triangle inequality separates $S$ and $Y$
    by at least $l$. The finite-support propagation bound gives
    $4|S|\|W\|\|B\|e^{v_\mu|t|-\mu l}$. Combine this with
    $\|[B,e^{iHt}We^{-iHt}]\|\le2\|W\|\|B\|$ and absorb the integer
    shift $2r+1$. At infinite anchor distance, let $l$ tend to infinity.
\end{proof}

\begin{lemma}[Gaussian truncation with weighted omitted interactions]
    \label{thm:hastings05_weighted_gaussian_truncation}
    \lean{ManyBody.norm_gaussianFilteredTerm_sub_truncated_le_anchor_sum}
    \leanok
    The retained family $s$ may be arbitrary. For every $\mu>0$,
    \begin{align}
        \|\mathcal F_{q,\Delta}^{H}(W)-\mathcal F_{q,\Delta}^{H_s}(W)\|
         & \le A_\mu(S)\|W\|T e^{v_\mu T}
        \sum_{j\notin s}\|h_j\|k_\mu(z,a_j)
        +2\|W\|e^{-(T\Delta)^2/(2q)}.
        \label{eq:hastings05_weighted_gaussian_truncation}
    \end{align}
    No cardinality bound is imposed on a region containing the retained
    interactions.
\end{lemma}
\begin{proof}\leanok
    \uses{thm:hastings05_anchor_lr, thm:hastings05_termwise_gaussian_truncation}
    The retained Hamiltonian satisfies the original local norm budget.
    Apply (\ref{eq:hastings05_anchor_lr}) to its evolution and each omitted
    $h_j$, with $b=a_j$ and $Y=X_j$. Bound $e^{v_\mu|t|}$ by
    $e^{v_\mu T}$ and use
    (\ref{eq:hastings05_termwise_gaussian_truncation}).
\end{proof}

\begin{lemma}[A bounded-degree tail controls truncation]
    \label{thm:hastings05_degree_gaussian_truncation}
    \lean{ManyBody.norm_gaussianFilteredTerm_sub_truncated_le_degree_tail}
    \leanok
    Suppose the graph degree is at most $D$, $B\ge0$ and
    $\sum_{j:\,a_j=x}\|h_j\|\le B$ at each anchor $x$.
    If $d_G(z,a_j)\ge m$ for every omitted $j$, choose
    $\mu>\log(D+1)$ and put $\beta=(\mu-\log(D+1))/2$.
    Then
    \begin{align}
        \|\mathcal F_{q,\Delta}^{H}(W)-\mathcal F_{q,\Delta}^{H_s}(W)\|
         & \le A_\mu(S)\|W\|T e^{v_\mu T}
        \frac{B e^{-\beta m}}{1-e^{-\beta}}
        +2\|W\|e^{-(T\Delta)^2/(2q)}.
        \label{eq:hastings05_degree_gaussian_truncation}
    \end{align}
\end{lemma}
\begin{proof}\leanok
    \uses{thm:hastings05_weighted_gaussian_truncation, thm:hastings05_degree_weighted_tail}
    Substitute the weighted tail estimate into
    (\ref{eq:hastings05_weighted_gaussian_truncation}).
\end{proof}

\begin{lemma}[Support and norm of the physical metric truncation]
    \label{thm:hastings05_metric_filter_properties}
    \lean{ManyBody.metric_truncated_filter_properties}
    \leanok
    If $\ell\ge R$, then $M_i(q,\Delta,\ell)$ is Hermitian, acts inside
    $B_\ell(i)$ and has norm at most $J$.
\end{lemma}
\begin{proof}\leanok
    \uses{thm:hastings05_metric_retained_support, thm:hastings_truncated_support, thm:hastings_gaussian_filtered_spectral}
    Every retained interaction and the initial $h_i$ act inside
    $B_\ell(i)$. Their evolution, and hence its Gaussian average, has the
    same support. Gaussian filtering preserves Hermitian operators and
    contracts the operator norm.
\end{proof}

\begin{lemma}[Uniform metric truncation and leakage budget]
    \label{thm:hastings05_metric_filter_budget}
    \lean{ManyBody.exists_metric_truncated_filter_budget}
    \leanok
    Fix $R>0$, $S\in\mathbb N$, $J\ge0$ and $\Delta>0$. There are
    $c_1,\lambda>0$ and $C\ge0$, chosen before $\Lambda$ and its
    Hamiltonian, such that every family satisfying the local metric
    assumptions and every $\ell\ge R$ obeys
    \begin{align}
        \|h_i\|e^{-q/2}
        +\|\mathcal F_{q,\Delta}^{H}(h_i)-M_i(q,\Delta,\ell)\|
         & \le\|h_i\|C e^{-\lambda\ell},
         & q                             & =c_1\ell\Delta.
        \label{eq:hastings05_metric_filter_budget}
    \end{align}
    One choice of constants is
    \begin{align}
        \mu    & =\log(S+1)+1,                                  & v                          & =4S^2Je^{4\mu},
               & c_1                                            & =\frac1{4R(v+1)},\nonumber                           \\
        \gamma & =\frac1{2R}-vc_1>0,                            &
        \delta & =\min(c_1\Delta,\gamma),                       & \lambda                    & =\frac\delta2,\nonumber \\
        A      & =\frac{4(S+1)e^{5\mu}Jc_1e^{1/2}}{1-e^{-1/2}}, &
        C      & =3+\frac{2A}{e\delta}.
        \label{eq:hastings05_metric_filter_constants}
    \end{align}
\end{lemma}
\begin{proof}\leanok
    \uses{thm:hastings05_degree_gaussian_truncation, thm:hastings05_metric_interaction_data, thm:hastings05_metric_graph_separation, lem:hastings05_gaussian_decay_absorption}
    The metric interaction graph has degree at most $S$, support diameter
    at most two and local norm budget $SJ$. Take anchors $a_j=j$, whose
    norm budget is $J$. Omitted anchors have graph distance at least
    $m=\lceil(\ell-R)/R\rceil$. With $\beta=1/2$,
    $e^{-\beta m}\le e^{1/2}e^{-\ell/(2R)}$.
    In (\ref{eq:hastings05_degree_gaussian_truncation}) take
    $T=c_1\ell$ and $q=c_1\ell\Delta$; its Gaussian tail becomes
    $2\|h_i\|e^{-q/2}$. Adding the spectral term gives
    $3\|h_i\|e^{-c_1\Delta\ell/2}+A\|h_i\|\ell e^{-\gamma\ell}$.
    The time-prefactor estimate bounds this by
    $\|h_i\|C e^{-\delta\ell/2}$.
\end{proof}

\begin{theorem}[Exponential localization of actual truncated terms]
    \label{thm:hastings05_metric_facld}
    \lean{ManyBody.exists_metric_truncated_filter_localization}
    \leanok
    Under the metric assumptions of
    Lemma~\ref{thm:hastings05_metric_filter_budget}, constants
    $c_1,\lambda>0$ and $C\ge0$, independent of the finite site set and
    Hamiltonian, give the following statements. Put $q=c_1\ell\Delta$
    and $M_i=M_i(q,\Delta,\ell)$ for $\ell\ge R$.
    Each $M_i$ is Hermitian, acts on $B_\ell(i)$ and has norm at most $J$,
    and
    \begin{align}
        \|\mathcal F_{q,\Delta}^{H}(h_i)-M_i\| & \le JC e^{-\lambda\ell}.
        \label{eq:hastings05_metric_filter_error}
    \end{align}
    If every eigenvalue of $H$ equals $E_0$ or is at least $E_0+\Delta$,
    let $P_0$ be the orthogonal projection onto the $E_0$-eigenspace.
    Every vector $\psi$ in that eigenspace satisfies
    \begin{align}
        \|(1-P_0)M_i\psi\| & \le JC e^{-\lambda\ell}\|\psi\|.
        \label{eq:hastings05_metric_facld}
    \end{align}
    In particular, if $\Omega$ is a normalized unique ground state with
    gap $\Delta$, then
    $\|M_i\Omega-\langle\Omega,M_i\Omega\rangle\Omega\|
        \le JC e^{-\lambda\ell}$.
    This is the exponential order asserted in (facld) of
    \cite[lines 322--348]{Hastings2006Solving}. The printed minimum of
    decay lengths is read as a minimum of decay rates, as explained in
    \cite{gap:hastings05_truncated_dynamics}.
\end{theorem}
\begin{proof}\leanok
    \uses{thm:hastings05_metric_filter_budget, thm:hastings05_metric_filter_properties, thm:hastings05_facld}
    The spectral leakage and approximation inequalities bound
    $\|(1-P_0)M_i\psi\|$ by the left side of
    (\ref{eq:hastings05_metric_filter_budget}) times $\|\psi\|$.
    Use $\|h_i\|\le J$. For a unique ground state,
    $P_0\phi=\langle\Omega,\phi\rangle\Omega$.
\end{proof}

\begin{theorem}[The sum of the physical metric approximations]
    \label{thm:hastings05_metric_filter_sum}
    \lean{ManyBody.exists_metric_truncated_filter_sum_bound}
    \leanok
    Fix $R,S,J,\Delta$ as in
    Lemma~\ref{thm:hastings05_metric_filter_budget}. Uniform constants
    $c_1,\lambda>0$ and $C\ge0$ give, for every admissible Hamiltonian and
    $\ell\ge R$,
    \begin{align}
        \Bigl\|\sum_iM_i(c_1\ell\Delta,\Delta,\ell)-H\Bigr\|
         & \le|\Lambda|JC e^{-\lambda\ell}.
        \label{eq:hastings05_metric_filter_sum}
    \end{align}
\end{theorem}
\begin{proof}\leanok
    \uses{thm:hastings05_metric_facld, thm:hastings05_sum_approximation}
    The full filtered terms sum to $H$. The finite-sum triangle
    inequality bounds the error by the sum of
    (\ref{eq:hastings05_metric_filter_error}).
\end{proof}


\section{The graph of a metric interaction range}
\label{sec:gaussian_filtered_metric_interaction_graph}

\begin{definition}[The metric interaction graph]
    \label{def:hastings05_metric_interaction_graph}
    \lean{ManyBody.metricRangeGraph, ManyBody.metricRangeBall, ManyBody.mem_metricRangeBall, ManyBody.coe_metricRangeBall, ManyBody.self_mem_metricRangeBall}
    \leanok
    Let $V$ be a finite metric space and $R\ge0$. Write $B_R(x)$ for its
    closed metric ball. Join distinct sites $x,y$ when $d(x,y)\le R$, obtaining
    a simple graph $G_R$. Every site belongs to its own interaction ball.
\end{definition}

\begin{lemma}[Local metric data give graph interaction bounds]
    \label{thm:hastings05_metric_interaction_data}
    \lean{ManyBody.metricRangeGraph_edist_le_one, ManyBody.metricRangeGraph_edist_le_two, ManyBody.metricRangeGraph_neighbor_card_le, ManyBody.metricRangeBall_graph_diameter_le_two, ManyBody.sum_norm_metricRangeBall_le, ManyBody.metricRangeBall_norm_budget_le}
    \leanok
    \uses{def:hastings05_metric_interaction_graph}
    Suppose $|B_R(x)|\le S$ for every site. The degree of $G_R$ is at most
    $S$, and every interaction ball has graph diameter at most two. If
    $\|h_i\|\le J$ with $J\ge0$, then
    $\sum_{i:\,x\in B_R(i)}\|h_i\|\le SJ$ at every site $x$.
\end{lemma}
\begin{proof}\leanok
    A graph neighbour lies in the metric ball. Two sites in one interaction
    ball are joined through its centre by at most two graph edges.
    By symmetry of the metric, the centres whose interaction balls contain
    $x$ lie in $B_R(x)$. Sum their norm bounds.
\end{proof}

\begin{lemma}[Physical separation and intrinsic graph distance]
    \label{thm:hastings05_metric_graph_separation}
    \lean{ManyBody.dist_le_range_mul_walkLength, ManyBody.dist_le_range_mul_graphDist, ManyBody.metricRangeGraph_edist_lower_ceil, ManyBody.metricRangeGraph_zero_edist_eq_top, ManyBody.metricRangeGraph_omitted_threshold}
    \leanok
    \uses{def:hastings05_metric_interaction_graph}
    Every graph walk of length $n$ has endpoints at physical distance at
    most $Rn$. Thus, on a connected component,
    $d(x,y)\le R d_{G_R}(x,y)$. For $R>0$, a site omitted by the source
    truncation $d(i,j)\le\ell-R$ has extended graph distance at least
    $\lceil(\ell-R)/R\rceil$ from $i$. At range zero, sites at positive
    physical distance lie in different graph components.
\end{lemma}
\begin{proof}\leanok
    Sum the edge-length bounds along a shortest walk. Graph distance is an
    integer, giving the ceiling bound; disconnected sites have infinite
    extended graph distance. At range zero, every reachable pair has
    physical distance zero.
\end{proof}

\begin{lemma}[The source's metric truncation preserves local support]
    \label{thm:hastings05_metric_retained_support}
    \lean{ManyBody.metricRangeBall_retained_support, ManyBody.metricRangeBall_initial_support}
    \leanok
    \uses{def:hastings05_metric_interaction_graph}
    If $d(i,j)\le\ell-R$, then $B_R(j)\subseteq B_\ell(i)$.
    If $\ell\ge R$, then $B_R(i)\subseteq B_\ell(i)$ as well.
\end{lemma}
\begin{proof}\leanok
    For $x\in B_R(j)$, the triangle inequality gives
    $d(i,x)\le d(i,j)+d(j,x)\le\ell-R+R=\ell$.
\end{proof}


\section{Connected sets in a graph of bounded degree}
\label{sec:gaussian_filtered_connected_set_count}

\begin{lemma}[Clusters through a vertex]
    \label{thm:hastings_connected_set_count_mem}
    \lean{ManyBody.card_connected_induce_mem_le}
    \leanok
    Let $G$ be a finite simple graph in which every vertex has degree at most
    $\Delta$, let $v$ be a vertex, and let $n$ be a natural number. At most
    $\Delta^{2(n-1)}$ sets of $n$ vertices contain $v$ and induce a connected
    subgraph of $G$, where $n-1$ is read as $0$ when $n=0$; no such set
    exists then.
\end{lemma}

\begin{proof}
    \leanok
    Let $S$ be such a set. Starting from $\{v\}$ and the walk of length zero,
    add the vertices of $S$ one at a time: as long as the set $T$ reached so
    far is a proper subset of $S$, a path inside $S$ from $v$ to a vertex of
    $S\setminus T$ leaves $T$ along an edge $wu$ with $w\in T$ and
    $u\in S\setminus T$. Inserting the detour $w\to u\to w$ into the current
    closed walk at a visit of $w$ adds two to its length and exactly the vertex
    $u$ to its vertex set. After $n-1$ steps this gives a closed walk from $v$
    of length $2(n-1)$ whose vertex set is $S$. The walk determines $S$, so
    distinct sets give distinct walks, and at most $\Delta^{L}$ walks of length
    $L$ start at $v$, since each step has at most $\Delta$ choices.
\end{proof}

\begin{lemma}[Clusters of a given size]
    \label{thm:hastings_connected_set_count}
    \lean{ManyBody.card_connected_induce_le, ManyBody.card_connected_induce_le_exponential}
    \leanok
    Let $G$ be a simple graph on $N$ vertices in which every vertex has degree
    at most $\Delta$, and let $n$ be a natural number. At most
    $N\Delta^{2(n-1)}$ sets of $n$ vertices induce a connected subgraph of $G$,
    with $n-1$ read as $0$ when $n=0$.
    In particular the number is at most $Nx^n$ for
    $x=\max(1,\Delta^2)$. The latter bound also holds for $n=0$, since the
    empty induced graph is not connected.
\end{lemma}

\begin{proof}
    \leanok
    \uses{thm:hastings_connected_set_count_mem}
    A set inducing a connected subgraph is nonempty, so it contains some
    vertex; sum the bound of Lemma~\ref{thm:hastings_connected_set_count_mem}
    over the $N$ vertices.
    For $n\ge1$, $\Delta^{2(n-1)}=(\Delta^2)^{n-1}\le x^{n-1}\le x^n$,
    since $x\ge1$. This choice also covers $\Delta=0$, when the connected
    singleton sets would prevent taking $x=0$.
\end{proof}


\section{Ball counts in integer lattices}
\label{sec:gaussian_filtered_cluster_truncation}

\begin{lemma}[Polynomial ball growth in the integer lattice]
    \label{lem:hastings_integer_lattice_ball_count}
    \lean{ManyBody.card_centres_integerLattice_ball_le, ManyBody.card_centres_integerLattice_ball_le_polynomial, ManyBody.card_centres_ball_le_polynomial_of_dist_le, ManyBody.card_centres_ball_le_polynomial_of_isometry}
    \leanok
    For distinct centres in $\mathbb Z^D$ with the coordinate maximum
    metric, every radius-$r$ ball contains at most
    $(2\lceil r\rceil+1)^D\le2^D(\lceil r\rceil+1)^D$ centres.
    The same bound holds for every metric on a finite integer-lattice
    domain that dominates coordinate maximum distance, including its
    intrinsic walk metric. More generally, it transfers through any
    injective coordinate map that does not increase distances.
    This is the polynomial volume growth, with $C=2^D$, used in the local
    index count of the cluster expansion.
\end{lemma}
\begin{proof}\leanok
    Every coordinate of a centre in the ball lies in an integer interval
    with $2\lceil r\rceil+1$ elements. Containment in the product of these
    intervals proves the first bound. The second follows from
    $2\lceil r\rceil+1\le2(\lceil r\rceil+1)$.
\end{proof}


\section{The first-moment Gaussian tail}
\label{sec:gaussian_filtered_gaussian_time_average}

\begin{lemma}[First absolute moment of a Gaussian tail]
    \label{thm:hastings_gaussian_absolute_moment_tail}
    \lean{GaussianFilter.integral_compl_Icc_abs_gaussianReal_le}
    \leanok
    For $h,T\ge0$ (with the right-hand side $0$ when $h=0$),
    \begin{align}
        \int_{|t|>T}|t|\,d\gamma_h(t)
        \le \sqrt{2h/\pi}\,e^{-T^2/(2h)}.
        \label{eq:hastings_gaussian_absolute_moment_tail}
    \end{align}
\end{lemma}
\begin{proof}\leanok
    For $h=0$, $\gamma_0=\delta_0$ gives $0$. For $h>0$, with
    $\varphi_h(t)=(2\pi h)^{-1/2}e^{-t^2/(2h)}$ and $(-h\varphi_h)'=t\varphi_h$,
    \begin{align}
        \int_{|t|>T}|t|\,d\gamma_h(t)=2\int_T^\infty t\varphi_h(t)\,dt
        =2h\varphi_h(T)=\sqrt{2h/\pi}\,e^{-T^2/(2h)}.
        \notag
    \end{align}
\end{proof}

\begin{lemma}[First absolute moment of a Gaussian]
    \label{thm:hastings_gaussian_absolute_moment}
    \lean{GaussianFilter.integral_abs_gaussianReal_le_sqrt_two_mul_div_pi, GaussianFilter.integral_abs_gaussianReal_le}
    \leanok
    For every $h\ge0$,
    $\int|t|\,d\gamma_h(t)\le\sqrt{2h/\pi}\le\sqrt h$.
\end{lemma}
\begin{proof}\leanok
    \uses{thm:hastings_gaussian_absolute_moment_tail}
    Set $T=0$ in~(\ref{eq:hastings_gaussian_absolute_moment_tail}) and use
    $2\le\pi$.
\end{proof}


\section{Powers against decaying exponentials}
\label{sec:gaussian_filtered_decay_slack}

\begin{lemma}[A power against an exponential]
    \label{thm:hastings_pow_exp}
    \lean{Real.pow_mul_exp_neg_le, Real.pow_mul_exp_neg_add_le}
    \leanok
    For $\delta>0$, a natural number $p$, a real number $r$ and $l\ge0$,
    \begin{align}
        l^pe^{-\delta l}    & \le\Bigl(\frac{p}{e\delta}\Bigr)^p,        &
        l^pe^{-(r+\delta)l} & \le\Bigl(\frac{p}{e\delta}\Bigr)^pe^{-rl}.
        \label{eq:hastings_pow_exp}
    \end{align}
\end{lemma}
\begin{proof}\leanok
    With $x=\delta l/p$ for $p\ge1$, $l^pe^{-\delta l}=(p/\delta)^p(xe^{-x})^p$ and
    $xe^{-x}\le e^{-1}$. The second bound follows from the first after
    splitting off $e^{-rl}$.
\end{proof}
